\documentclass[10pt,a4paper]{amsart}
\usepackage[margin=19mm]{geometry}
\usepackage{amsmath,amssymb,mathtools}
\usepackage[scr=rsfs,cal=euler]{mathalfa}
\usepackage{xfrac}
\usepackage[unicode,hidelinks,bookmarks=false]{hyperref}
\newcommand{\ie}{i.e.\ }
\newcommand{\cf}{cf.\ }
\newcommand{\resp}{respectively\ }
\newcommand{\Subsection}{\subsection}
\providecommand{\nobreakdash}{\nobreak\hbox{-}}
\setlength{\emergencystretch}{2em}
\multlinegap0pt
\usepackage{dsfont}
\usepackage[shortlabels]{enumitem}
\usepackage{float}
\usepackage{amssymb}
\usepackage{multirow}
\usepackage{xcolor}
\usepackage{cancel}
\usepackage{bm}
\usepackage{mathtools}
\newtheorem{thm}{Theorem}
\newcommand{\ba}{\bm{a}}
\newcommand{\bx}{\bm{x}}
\title{Linear minimum-variance approximants for noisy data}
\author{Sergio L\'opez-Ure\~na, Dionisio F. Y\'a\~nez}
\date{Source: arXiv:2412.01287v4 (CC BY 4.0)}
\begin{document}
\maketitle

\noindent\emph{Source and licence.} Linear minimum-variance approximants for noisy data. Version arXiv:2412.01287v4 (CC BY 4.0), distributed by arXiv under the Creative Commons Attribution 4.0 International licence, http://creativecommons.org/licenses/by/4.0/, as recorded in the arXiv licence metadata for this exact version and shown on the abstract page https://arxiv.org/abs/2412.01287v4. Full licence notice: \texttt{licenses/arxiv-2412.01287-CC-BY-4.0.txt}.\par\smallskip

\noindent\textit{Editorial scope.} Counted unit: the article's thm thm:orthonormal (source0.tex lines 333--339). The article's complete original proof of it is delivered with it (source0.tex lines 340--353, one \texttt{proof} environment of 14 lines). No mathematics is written, repaired or re-derived: the statement and the proof are the article's own lines, in the article's own notation, and no retained line is substituted or re-flowed.  Retained uncounted as the source-local context the proof needs: lines 315--317; lines 325--327.  Nothing was omitted from any retained span, so the package carries no omission record.  No retained line contains a citation, so the wrapper carries no bibliography and no bibliography entry.  No retained line contains a figure environment, an image-inclusion command or drawing code, so no image is delivered and the package declares no asset dependency.  Editorial formatting only, in addition: an AMS \texttt{amsart} preamble that restates the article's own declarations and macros used by the retained lines, namely the environments \texttt{thm} on separate counters, together with the standard packages those lines need.  The retained environments print in this document as Theorem 1 in order of appearance, whereas the article prints the same items with the numbering of its own full document.  The complete native source of the article as distributed in the arXiv e-print for this exact version is bundled unchanged as \texttt{sources/169448/source0.tex}.
\begin{thm} \label{thm:caracterization_of_optimal}
	Let $\Psi$ be an approximant with coefficients $\ba$ that reproduces $\Pi_{d-1}$. It is minimum-variance if and only if $\exists Q\in\Pi_{d-1}$ such that $\ba = \hat \Omega^{-1} Q|_{\bx}$.
\end{thm}

\begin{equation} \label{eq:d_times_d}
    \sum_{i=1}^{N} (\hat \Omega^{-1} Q|_{\bx})_i x_i^s = t_0^s, \qquad s = 0,1,\ldots,d-1.
\end{equation}

\begin{thm}\label{thm:orthonormal}
    Consider the inner product $(P,Q) = P|_{\bx}^T \hat\Omega^{-1} Q|_{\bx}$ for the space $\Pi_{N-1}$.
    Let $\{P^s\}_{s=0}^{d-1}$ be a family of orthonormal polynomials such that $P^s\in\Pi_s$. The coefficients of the minimum-variance approximant for $\hat\Omega$, $t_0$, $\bx$ and $d-1$ are
    \[
        \ba = \hat\Omega^{-1} Q|_{\bx}, \quad Q:=\sum_{j=0}^{d-1} P^j(t_0) P^j.
    \]
\end{thm}

\begin{proof}
    By applying Theorem \ref{thm:caracterization_of_optimal}, it remains to show that with such $Q$ the approximant certainly reproduces $\Pi_{d-1}$. That is, we may show that the approximant is exact for a basis of $\Pi_{d-1}$, as done in \eqref{eq:d_times_d}. In particular, for the basis $\{P^s\}_{s=0}^{d-1}$, we may check that
    \[
    \sum_{i=1}^{N} (\hat \Omega^{-1} Q|_{\bx})_i P^s(x_i) = P^s(t_0), \qquad s = 0,1,\ldots,d-1.
    \]
    Indeed, for $Q=\sum_{j=0}^{d-1} P^j(t_0) P^j$,
    \begin{align*}
        \sum_{i=1}^{N} (\hat \Omega^{-1} Q|_{\bx})_i P^s(x_i)
        = (Q|_{\bx})^T \hat \Omega^{-1} P^s|_{\bx}
        = \left(Q,P^s\right)
        = \sum_{j=0}^{d-1} P^j(t_0) \left(P^j,P^s\right)
        = P^s(t_0).
    \end{align*}
\end{proof}

\end{document}
